% Options for packages loaded elsewhere
\PassOptionsToPackage{unicode}{hyperref}
\PassOptionsToPackage{hyphens}{url}
\PassOptionsToPackage{dvipsnames,svgnames,x11names}{xcolor}
\documentclass[
  a4paper,
]{ltjsarticle}
\usepackage{xcolor}
\usepackage[margin=25mm]{geometry}
\usepackage{amsmath,amssymb}
\setcounter{secnumdepth}{-\maxdimen} % remove section numbering
\usepackage{iftex}
\ifPDFTeX
  \usepackage[T1]{fontenc}
  \usepackage[utf8]{inputenc}
  \usepackage{textcomp} % provide euro and other symbols
\else % if luatex or xetex
  \usepackage{unicode-math} % this also loads fontspec
  \defaultfontfeatures{Scale=MatchLowercase}
  \defaultfontfeatures[\rmfamily]{Ligatures=TeX,Scale=1}
\fi
\usepackage{lmodern}
\ifPDFTeX\else
  % xetex/luatex font selection
  \setmainfont[]{DejaVu Serif}
  \setsansfont[]{DejaVu Sans}
  \setmonofont[]{DejaVu Sans Mono}
\fi
% Use upquote if available, for straight quotes in verbatim environments
\IfFileExists{upquote.sty}{\usepackage{upquote}}{}
\IfFileExists{microtype.sty}{% use microtype if available
  \usepackage[]{microtype}
  \UseMicrotypeSet[protrusion]{basicmath} % disable protrusion for tt fonts
}{}
\makeatletter
\@ifundefined{KOMAClassName}{% if non-KOMA class
  \IfFileExists{parskip.sty}{%
    \usepackage{parskip}
  }{% else
    \setlength{\parindent}{0pt}
    \setlength{\parskip}{6pt plus 2pt minus 1pt}}
}{% if KOMA class
  \KOMAoptions{parskip=half}}
\makeatother
\usepackage{longtable,booktabs,array}
\newcounter{none} % for unnumbered tables
\usepackage{calc} % for calculating minipage widths
% Correct order of tables after \paragraph or \subparagraph
\usepackage{etoolbox}
\makeatletter
\patchcmd\longtable{\par}{\if@noskipsec\mbox{}\fi\par}{}{}
\makeatother
% Allow footnotes in longtable head/foot
\IfFileExists{footnotehyper.sty}{\usepackage{footnotehyper}}{\usepackage{footnote}}
\makesavenoteenv{longtable}
\setlength{\emergencystretch}{3em} % prevent overfull lines
\providecommand{\tightlist}{%
  \setlength{\itemsep}{0pt}\setlength{\parskip}{0pt}}
\usepackage{xurl}
\usepackage{bookmark}
\IfFileExists{xurl.sty}{\usepackage{xurl}}{} % add URL line breaks if available
\urlstyle{same}
\hypersetup{
  pdftitle={有界なヒルベルト空間と八元数による時空・素粒子の創発},
  pdfauthor={木原 範昭（株式会社WFシステム）　ORCID 0009-0004-6753-4020},
  colorlinks=true,
  linkcolor={Maroon},
  filecolor={Maroon},
  citecolor={Blue},
  urlcolor={Blue},
  pdfcreator={LaTeX via pandoc}}

\title{有界なヒルベルト空間と八元数による時空・素粒子の創発}
\usepackage{etoolbox}
\makeatletter
\providecommand{\subtitle}[1]{% add subtitle to \maketitle
  \apptocmd{\@title}{\par {\large #1 \par}}{}{}
}
\makeatother
\subtitle{―
光子の幾何学・中心投影から、背景時空を仮定しない四元数・八元数動力学の課題整理まで
―}
\author{木原 範昭（株式会社WFシステム）　ORCID 0009-0004-6753-4020}
\date{2026年11月16日　Version DOI 10.5281/zenodo.24267348}

\ltjsetparameter{jacharrange={-2,-3}}
\begin{document}
\maketitle

\textbf{著者：} 木原範昭（株式会社WFシステム）　ORCID
0009-0004-6753-4020\\
\textbf{日付：} 2026-11-16　　\textbf{版：} v1.0\\
\textbf{Version DOI：} 10.5281/zenodo.24267348　　\textbf{Concept DOI：}
10.5281/zenodo.24267347\\
\textbf{ライセンス：} CC BY 4.0

\begin{center}\rule{0.5\linewidth}{0.5pt}\end{center}

\subsection{要旨}\label{ux8981ux65e8}

本稿は、光子の幾何学的性質（四次元時空における静止）の直観から出発し、中心投影モデル、四元数による位相勾配の保持、そして八元数の非結合性を用いた「外部背景時空を仮定しないヒルベルト空間からの粒子と相互作用の自己組織化」に至るまでの思考実験と対話の記録である。

\begin{center}\rule{0.5\linewidth}{0.5pt}\end{center}

\subsection{はじめに：本稿の位置づけと既存研究との接続}\label{ux306fux3058ux3081ux306bux672cux7a3fux306eux4f4dux7f6eux3065ux3051ux3068ux65e2ux5b58ux7814ux7a76ux3068ux306eux63a5ux7d9a}

本稿は、あらかじめ4次元の背景時空を舞台として置かず、有界・有限なヒルベルト空間の状態関係と四元数・八元数の代数構造のみから、時空・粒子・相互作用を再構成しようとする思考実験の記録である。

このような出発点は一見すると標準的な教科書的記述から大きく離れて見えるが、問題意識と数学的道具立ての双方において、現代の理論物理学で研究されている複数の潮流と接続している。

\begin{enumerate}
\def\labelenumi{\arabic{enumi}.}
\tightlist
\item
  \textbf{創発的時空（Emergent Spacetime）と有界な状態空間：}
  4次元多様体を前提とせず、量子状態間の相関や部分系への分解から空間次元・幾何学を再構成する研究（Quantum
  Mereology
  やテンソルネットワークによる時空創発）、およびド・ジッター時空の事象の地平面に対応づけて有限次元・有界なヒルベルト空間を考察する議論と問題意識を共有している。
\item
  \textbf{スピノル・多元体代数による時空と素粒子の記述：}
  複素数ペア（スピノル）とホップ写像 {[}X5{]}{[}X6{]}
  を介して空間座標や光線を捉える視点はツイスター理論や量子情報のブロッホ球表現と構造を同じくし、八元数の非結合性・トライアリティ・極小左イデアルを用いて粒子種や対称性を記述する試みは、多元体による素粒子模型の代数的研究
  {[}X8{]}{[}X9{]} と接点を持つ。
\end{enumerate}

なお、本稿は完成した理論の主張ではなく、現在進行中の思考整理である。既存の数学的恒等式・物理的解釈と、本稿固有の作業仮説、および未解決の導出課題（内部位相を書き換える具体的更新則
{[}S9{]}--{[}S11{]} と、\(N^2\)
全結合からの因子分解の数学的証明）の境界については、末尾の「区分」に明示した通り厳密に切り分けて位置づけている。

{\def\LTcaptype{none} % do not increment counter
\begin{longtable}[]{@{}
  >{\raggedright\arraybackslash}p{(\linewidth - 6\tabcolsep) * \real{0.2500}}
  >{\raggedright\arraybackslash}p{(\linewidth - 6\tabcolsep) * \real{0.2500}}
  >{\raggedright\arraybackslash}p{(\linewidth - 6\tabcolsep) * \real{0.2500}}
  >{\raggedright\arraybackslash}p{(\linewidth - 6\tabcolsep) * \real{0.2500}}@{}}
\toprule\noalign{}
\begin{minipage}[b]{\linewidth}\raggedright
段階
\end{minipage} & \begin{minipage}[b]{\linewidth}\raggedright
数学的構造・数体系
\end{minipage} & \begin{minipage}[b]{\linewidth}\raggedright
幾何学的写像
\end{minipage} & \begin{minipage}[b]{\linewidth}\raggedright
物理的対応（作業仮説）
\end{minipage} \\
\midrule\noalign{}
\endhead
\bottomrule\noalign{}
\endlastfoot
第1段階 & ミンコフスキー計量 \((x,y,z,t,\tau)\) & 5次元光円錐の断面 &
光子の固有時ゼロ（\(ds=0\)）・物質の時間進行 \\
第2段階 & 虚時間 \(R^2 = t^2 + (i\tau)^2\) & 中心投影（ベルトラミ座標）
& ド・ジッター時空・事象の地平面 \(R\) と重力歪み \\
第3段階 & 四元数 \(\mathbb{H}\)（共役な2複素数） & ホップ写像
\(S^1 \hookrightarrow S^3 \to S^2\) &
外部座標を借りない3方向位相勾配の保持（光子） \\
第4段階 & 八元数 \(\mathbb{O}\)（共役な2四元数） & ホップ写像
\(S^3 \hookrightarrow S^7 \to S^4\) & 閉じた4次元時空 \(S^4\)
と内部3自由度・凝縮体の自己組織化 \\
\end{longtable}
}

\begin{center}\rule{0.5\linewidth}{0.5pt}\end{center}

\subsection{第1段階：光子と四次元時空の幾何学}\label{ux7b2c1ux6bb5ux968eux5149ux5b50ux3068ux56dbux6b21ux5143ux6642ux7a7aux306eux5e7eux4f55ux5b66}

\textbf{【木原】} 光子が常に光速で飛ぶように見えるのはなぜですか？

\textbf{【Gemini】}
光子が常に光速（秒速約30万km）で飛ぶ理由は、光子の静止質量がゼロであるためです。物理学の各視点からは次のように整理されます。

\begin{enumerate}
\def\labelenumi{\arabic{enumi}.}
\tightlist
\item
  \textbf{相対性理論の視点：} 粒子のエネルギー \(E\)、運動量
  \(p\)、静止質量 \(m_0\) の関係式 \(E^2 = (pc)^2 + (m_0 c^2)^2\)
  において、\(m_0 > 0\) の物質は静止していても \(E = m_0 c^2\)
  のエネルギーを持ち、光速 \(c\)
  に近づけるとエネルギーが発散するため到達できません。一方、光子では
  \(m_0 = 0\) より \(E = pc\) となり、エネルギーを持つためには必ず運動量
  \(p\) を持って光速 \(c\) で運動し続ける必要があります（静止すると
  \(E = 0\) となります）。
\item
  \textbf{素粒子物理学の視点：}
  空間を満たすヒッグス場と相互作用する粒子（電子・クォーク等）は質量を獲得しますが、光子はヒッグス場と相互作用しないため、因果律の上限速度である光速で伝播します。
\item
  \textbf{電磁気学の視点：}
  変化する電場と磁場の連鎖（電磁波）が真空中を伝わる速さは、真空の誘電率
  \(\varepsilon_0\) と透磁率 \(\mu_0\) のみによって
  \(c = 1/\sqrt{\varepsilon_0 \mu_0}\) と定まります。
\item
  \textbf{補足（媒質中の伝播）：}
  水やガラス中で見かけの速度が遅くなる場合でも、光子自体は分子間の真空を常に光速で飛んでおり、物質内の電子との電磁的相互作用（吸収・再放出や干渉）により通過に時間がかかるためです。
\end{enumerate}

\textbf{【木原】}
光子だけが四次元時空で静止していて、陽子や電子だけが時間の速度で移動していると解釈してもよいのでは？

\textbf{【Gemini】}
その視点は、特殊相対性理論を「4次元時空の幾何学（ミンコフスキー時空）」として捉えた記述に対応します。前回の説明が「空間の中で粒子がどう振る舞うか（エネルギーと質量の視点）」であったのに対し、「時空全体の構造の中でどう存在しているか（幾何学の視点）」から同じ現象を記述したものです。

\begin{enumerate}
\def\labelenumi{\arabic{enumi}.}
\tightlist
\item
  \textbf{陽子や電子は時間軸を光速で移動している：}
  4次元時空の幾何学的解釈では、すべての物体が4次元時空内を一定の速さ
  \(c\)
  で進んでいるとみなせます。空間的に静止した陽子や電子は移動能力の100\%を時間軸方向（速度
  \(c\)）に使い、空間移動が増えるとその分だけ時間軸方向の進行が削られます（時間の遅れ）。光子は移動の100\%を空間方向に振り切っているため、時間軸方向の進行がゼロになります。
\item
  \textbf{光子にとっての時空間隔はゼロ：} 4次元時空における時空間隔
  \(ds\) は \[ds^2 = -(c\,dt)^2 + dx^2 + dy^2 + dz^2\]
  で表されます。光子は時間 \(dt\) の間に空間距離
  \(\sqrt{dx^2 + dy^2 + dz^2} = c\,dt\) を進むため、常に \(ds = 0\)
  となり、軌跡に沿った固有時間は経過しません。
\item
  \textbf{表現の整理：}
  厳密には「4次元時空の座標そのもので一点に止まっている」のではなく、ミンコフスキー図上では斜め45度の光円錐に沿って座標
  \((x, t)\)
  を変化させつつ、「その軌跡に沿った固有時間の経過（4次元的な距離）がゼロである」という表現に対応します。
\end{enumerate}

\textbf{【木原】} 光円錐の写像ですね。\(x^2+y^2+z^2=t^2+R^2\)

\textbf{【Gemini】} その式（光速 \(c = 1\)
とした自然単位系）は、4次元時空の現象を「より高次元の光円錐（ヌル構造）からの断面（写像）」として捉える幾何学に対応します。

\begin{enumerate}
\def\labelenumi{\arabic{enumi}.}
\tightlist
\item
  \textbf{高次元光円錐の断面（写像）：} 3次元空間距離の2乗
  \(x^2+y^2+z^2\) が、時間 \(t\) とパラメータ \(R\)
  の2乗和と釣り合っており、次元を1つ上げた5次元空間の光円錐を特定の値
  \(R\) でスライスして4次元時空へ写像した形と解釈できます。
\item
  \textbf{不変双曲面：} 移項して \(x^2+y^2+z^2-t^2=R^2\)
  とすると、\(R=0\) で光円錐そのもの、\(R>0\)
  で光円錐の外側に張り付く一葉双曲面となります。
\item
  \textbf{補足（「時間の速度」との対応）：}
  物質が時間軸を進む直観をピタゴラスの定理に対応させる場合は、物質の固有時
  \(\tau\) を用いて \(t^2 = (x^2+y^2+z^2) + \tau^2\)
  と書けます。座標時間 \(t\) を斜辺、空間距離を底辺、固有時 \(\tau\)
  を高さとする直角三角形（エプスタインの幾何学）において、光子（\(\tau=0\)）は光円錐の境界に、陽子や電子（\(\tau>0\)）は内側に位置します。
\end{enumerate}

\textbf{【木原】} 固有時間 \(\tau\) を使った表現
\(x^2+y^2+z^2=t^2-\tau^2\) の方が標準理論に対応させやすいですね。

\textbf{【Gemini】} \((x^2 + y^2 + z^2) + \tau^2 = t^2\)
は特殊相対性理論の標準的な計量に対応し、光円錐の写像と固有時間 \(\tau\)
を結びつける式になります。

\begin{enumerate}
\def\labelenumi{\arabic{enumi}.}
\tightlist
\item
  \textbf{座標時間 \(t\) を斜辺とする関係式：} 光子（\(\tau = 0\)）は
  \(x^2+y^2+z^2=t^2\)
  として光円錐の境界に位置し、陽子や電子（\(\tau > 0\)）は固有時間
  \(\tau\) の分だけ空間距離の2乗が \(t^2\)
  から差し引かれ、光円錐の内側に閉じ込められます（光速を超えられない幾何学的理由）。
\item
  \textbf{5次元光円錐からの写像としての解釈：}
  \((x^2+y^2+z^2) + \tau^2 = t^2\) は、3次元空間 \((x, y, z)\)
  に第4の空間的成分 \(\tau\)
  を加えた5次元的な光円錐の式と同型です。すべての存在がこの光円錐上を光速で走り、速度ベクトルを空間
  \((x,y,z)\) のみに向けるものが光子、一部または全部を固有時間 \(\tau\)
  に向けるものが質量を持つ粒子として射影される、と統一的に表せます。
\item
  \textbf{量子力学との接点：}
  場の量子論（ディラック方程式の性質）においても、質量を持つ粒子はコンプトン波長以下のミクロなスケールで光速のジグザグ運動（ツィッターベーグング）を行い、その平均速度が光速未満に見えるという記述が現れます。
\end{enumerate}

\begin{center}\rule{0.5\linewidth}{0.5pt}\end{center}

\subsection{第2段階：中心投影とド・ジッター時空}\label{ux7b2c2ux6bb5ux968eux4e2dux5fc3ux6295ux5f71ux3068ux30c9ux30b8ux30c3ux30bfux30fcux6642ux7a7a}

\textbf{【木原】} ここで \(R^2=t^2+(i\tau)^2\) として表現すると中心投影
\[x^2+y^2+z^2=R^2\] になります。局所曲率半径 \(R\)
を使った重力歪みのある場合に見えます。

\textbf{【Gemini】}
\[t^2 - \tau^2 = t^2 + (i\tau)^2 = R^2 \implies x^2 + y^2 + z^2 = R^2\]
という変換は、特殊相対性理論（平坦な時空）から一般相対性理論（曲がった時空）へ接続する幾何学的手続きに対応します。

\begin{enumerate}
\def\labelenumi{\arabic{enumi}.}
\tightlist
\item
  \textbf{虚時間 \((i\tau)\) による双曲幾何から球面幾何への転換：}
  固有時間に虚数単位 \(i\) を掛けて \((i\tau)^2 = -\tau^2\)
  とすること（ウィック回転／虚時間）により、双曲的な関係が半径 \(R\)
  の球面方程式 \(x^2+y^2+z^2=R^2\) として扱えるようになります。
\item
  \textbf{中心投影と局所曲率半径 \(R\) による重力の記述：}
  \(x^2+y^2+z^2=R^2\)
  を高次元空間からの中心投影（グノモン投影／ベルトラミ・クラインモデル
  {[}X1{]}）として見ると、球面上の測地線が投影先の平坦空間では直線として写像されます
  {[}S1{]}。任意の重力場であっても、局所慣性系を切り出せば「局所曲率半径
  \(R\) を持つ球面の接空間への中心投影」として近似する視点と対応します。
\item
  \textbf{光子と物質の対比：} \(R^2 = t^2 + (i\tau)^2\)
  より、光子（\(\tau=0\)）では \(R=t\)
  となり光円錐（ホライズン）の境界を走ります。一方、質量を持つ物質（\(\tau>0\)）では
  \((i\tau)^2 = -\tau^2\) により実効的な半径 \(R\) が \(t\)
  より縮み（\(R<t\)）、固有時 \(\tau\) を持つことが局所曲率半径 \(R\)
  を縮めて重力に対応する歪みを生む、という解釈になります。
\end{enumerate}

\textbf{【木原】}
これは今年4月から着想していたモデルで、ド・ジッター時空 {[}X2{]}
という非等方にも見えるモデルを中心投影として理解できるモデルです
{[}S1{]}{[}S2{]}。

\textbf{【Gemini】} 宇宙項によって加速膨張するド・ジッター時空 {[}X2{]}
を、FLRW
計量ではなく高次元の超曲面からの中心投影（射影座標／ベルトラミ座標）として捉え直すと、以下の構造が明確になります
{[}S1{]}{[}S2{]}。

\begin{enumerate}
\def\labelenumi{\arabic{enumi}.}
\tightlist
\item
  \textbf{測地線の直線化：}
  ド・ジッター時空上のすべての測地線（光の軌道および粒子の慣性運動）が、投影先の空間では直線として写像されます。
\item
  \textbf{事象の地平面の可視化：}
  因果律の限界である宇宙の地平面（ハッブル半径）が、そのまま
  \(x^2+y^2+z^2=R^2\) の球面境界として表されます。
\item
  \textbf{非定常性・非等方性の還元：}
  観測者の視点によって空間の伸び方や時間経過が異なって見える性質も、5次元空間における対称な双曲面（あるいは虚時間
  \(i\tau\)
  を導入した4次元球面）をどの角度・断面から中心投影しているかという射影の違いとして整理できます。
\end{enumerate}

\begin{center}\rule{0.5\linewidth}{0.5pt}\end{center}

\subsection{第3段階：複素数から四元数・八元数への拡張}\label{ux7b2c3ux6bb5ux968eux8907ux7d20ux6570ux304bux3089ux56dbux5143ux6570ux516bux5143ux6570ux3078ux306eux62e1ux5f35}

\textbf{【木原】}
このモデルでの光子を複素数波で表現せず、四元数で表現すると分かりやすい写像になりませんか？

\textbf{【Gemini】}
中心投影モデルにおいて、光子を複素数波（\(\mathbb{C}\)）ではなく四元数（\(\mathbb{H}\)）で表現すると、「時空の幾何学（中心投影）」と「波の位相・偏光」を1つの写像として統合できます。

\begin{enumerate}
\def\labelenumi{\arabic{enumi}.}
\tightlist
\item
  \textbf{空間3軸の虚数単位 \((\mathbf{i}, \mathbf{j}, \mathbf{k})\)
  への吸収：} 通常の複素数波
  \(e^{i(\mathbf{k}\cdot\mathbf{x}-\omega t)}\)
  では虚数単位が1つのため、3次元空間の進行方向や偏光を外から手動で与える必要があります。四元数
  \(q = t + x\mathbf{i} + y\mathbf{j} + z\mathbf{k}\)（\(\mathbf{i}^2=\mathbf{j}^2=\mathbf{k}^2=\mathbf{i}\mathbf{j}\mathbf{k}=-1\)）を用いると、空間ベクトル
  \(\mathbf{r} = x\mathbf{i} + y\mathbf{j} + z\mathbf{k}\) について
  \(\mathbf{r}^2 = -(x^2+y^2+z^2)\)
  となり、\(t^2 + \mathbf{r}^2 = t^2 - (x^2+y^2+z^2) = \tau^2\)
  というミンコフスキー計量の符号反転が代数構造に内蔵されます。
\item
  \textbf{球面 \(S^3\) 上の回転写像：} 単位純虚四元数
  \(\mathbf{n}\)（\(\mathbf{n}^2=-1\)）に対するオイラーの公式
  \(e^{\mathbf{n}\theta} = \cos\theta + \mathbf{n}\sin\theta\)
  により、大きさ一定の四元数 \(|q|=R\) の集合は3次元超球面 \(S^3\)
  を形成します。光子（\(\tau=0\)）の伝播は
  \(q(\theta) = q_0 e^{\mathbf{n}\theta}\) という \(S^3\)
  上の大円（測地線）に沿った回転写像となり、中心投影により3次元空間を進行軸
  \(\mathbf{n}\)
  の周りに回転しながら進むらせん波（円偏光）として表されます。
\item
  \textbf{偏光（スピン）の幾何学化：}
  四元数では積の順序（左から掛けるか右から掛けるか）と回転軸
  \(\mathbf{n}\) の向きによってヘリシティ（スピン
  \(\pm 1\)）が表され、\(S^3\)
  から3次元空間への射影（ホップ写像）により位相ファイバーが円環構造（ホップリンク）として写像されます。
\item
  \textbf{マクスウェル方程式との接続：} 四元数微分演算子
  \(D = \partial_t + \mathbf{i}\partial_x + \mathbf{j}\partial_y + \mathbf{k}\partial_z\)
  を電磁ポテンシャル四元数に作用させると、真空中のマクスウェル方程式全体が（複素四元数表現において）\(DF = 0\)
  の1式に凝縮され、四元数正則性と光円錐上の写像が対応します。
\end{enumerate}

\textbf{【木原】} 四元数を知ったのは僅か1カ月前ですが
{[}S3{]}、100年以上前にマクスウェルはそちらから定式化に入っていたのですね
{[}X3{]}。

\textbf{【Gemini】} 1873年にマクスウェルが電磁気学をまとめた際
{[}X3{]}、ハミルトンの四元数を用いて記述していましたが、後にヘヴィサイドやギブズらによってスカラーと3次元ベクトル（内積・外積）へ分離されたため、現代の物理教育では四元数とは切り離して教えられる経緯があります。

\textbf{【木原】}
複素数の状態とその干渉を相互作用として考えても、3次元座標を読み出す自由度が足らず、光子が自身の座標（三方向の位相勾配）を保持（レジスト）できる写像を作れませんでした。そのため最初は共役な2つの複素数で光子を表現する方向で考えていました
{[}S3{]}{[}S4{]}{[}S5{]}。

\textbf{【Gemini】}
「共役な2つの複素数」による構成は、数学のケイリー＝ディクソン構成
{[}X4{]}（複素数から四元数を生成する手続き）そのものに対応します
{[}S3{]}。

\begin{enumerate}
\def\labelenumi{\arabic{enumi}.}
\tightlist
\item
  \textbf{複素数1つで3次元座標を保持できない理由：} 複素波
  \(\psi = A e^{i\theta}\) は実2自由度しか持たず、位相 \(\theta\)
  は1次元円周 \(S^1\) 上に留まるため、3次元空間の進行方向 \((x, y, z)\)
  の位相勾配を波の内部に記録できず、外部の座標グリッドを借りる必要が生じます。
\item
  \textbf{共役な2つの複素数から四元数への移行：} 2つの複素数
  \((z_1, z_2) \in \mathbb{C}^2\) を並べて \(q = z_1 + z_2 \mathbf{j}\)
  とおくと、積
  \(qw = (z_1 w_1 - z_2 \bar{w}_2) + (z_1 w_2 + z_2 \bar{w}_1)\mathbf{j}\)
  に複素共役が現れ、一意に四元数 \(\mathbb{H}\)
  の代数（ノルム保存）が得られます。
\item
  \textbf{ホップ写像 {[}X5{]} による3次元座標の生成：} \((z_1, z_2)\)
  とその共役 \((\bar{z}_1, \bar{z}_2)\) から、外部座標を借りずに
  \[x = 2\,\mathrm{Re}(z_1 \bar{z}_2),\quad y = 2\,\mathrm{Im}(z_1 \bar{z}_2),\quad z = |z_1|^2 - |z_2|^2\]
  を構成すると、
  \[x^2 + y^2 + z^2 = (|z_1|^2 + |z_2|^2)^2 = R^2 \quad (\text{ただし } |z_1|^2 + |z_2|^2 = R)\]
  となり、中心投影の球面方程式が直接得られます。全体の大きさ \(R\)
  を固定した残りの3自由度（\(S^3\)）が、3次元空間の方向（\(S^2\)
  上の2自由度）と偏光位相（\(S^1\)
  の1自由度）に分解されます（ツイスター理論のようにスピノルから時空を創発させる問題意識とも重なります）。
\end{enumerate}

\textbf{【木原】} 既存の数学的構成（ケイリー＝ディクソン構成
{[}X4{]}）を、物理的要請から逆にたどる形で到達したことになりますね。

\textbf{【Gemini】}
完成された代数形式を先に当てはめるのではなく、「外部座標を借りずに3方向の位相勾配を保持できる最小自由度」という物理的要請から共役な2複素数・四元数へ至ったため、各自由度の幾何学的役割と直結した形で整理されています。

\textbf{【木原】} ここまでは1カ月前に到達しています
{[}S3{]}{[}S4{]}。さらにこれでは質量と電荷を分離して表す自由度が足りないことも半年前から気づいています
{[}S6{]}{[}S7{]}。つまり \(x, y, z, t, \tau\)
では3自由度程度不足しており、\(x, y, z, t, q_1, q_2, q_3, \tau\)
という共役な四元数2つの八元数が電荷を持つ粒子に対応するのではないか、と1カ月前に考えました
{[}S5{]}{[}S8{]}。

\textbf{【Gemini】} ケイリー＝ディクソン構成 {[}X4{]}
をもう一段階進め、2つの四元数 \(a, b \in \mathbb{H}\)
とその共役から構成される \(O = a + b\ell\)
が8次元の八元数（\(\mathbb{O}\)）です。

\begin{enumerate}
\def\labelenumi{\arabic{enumi}.}
\tightlist
\item
  \textbf{ノルム保存の上限としての8次元（フルヴィッツの定理
  {[}X7{]}）：} 中心投影の幾何学を保つには積のノルム保存
  \(|AB| = |A||B|\) が不可欠ですが、フルヴィッツの定理 {[}X7{]}
  によりこれを満たす実ノルム多元体は 1, 2, 4, 8
  次元のみに限られます。16次元（十六元数）では零因子が現れてノルム保存が破綻するため、8次元が上限になります。
\item
  \textbf{外部時空ブロックと内部自由度ブロックの分離：} 8自由度を
  \((t, x, y, z)\) と \((\tau, q_1, q_2, q_3)\)
  の2ブロックに分ける見方は、八元数を用いた標準模型の代数的研究 {[}X9{]}
  の構成とも接点を持ちます。\(\tau\)（1自由度）が固有時間・質量に対応し、残る3つの虚軸
  \((q_1, q_2, q_3)\) が電荷・色荷（\(U(1)\) や
  \(SU(3) \subset G_2\)）に対応する内部自由度の候補となります（具体的なゲージ群との同定は接続課題です）{[}S5{]}{[}S8{]}。
\item
  \textbf{非結合性の出現：}
  四元数で非可換性（\(ab \neq ba\)）が現れたように、八元数では結合法則が破れます（\((ab)c \neq a(bc)\)）。結合法則を満たす四元数（光子）が自由に伝播できるのに対し、八元数では内部3自由度
  \((q_1, q_2, q_3)\)
  の非結合的な絡み合いが、単体で取り出せない閉じ込めや3世代構造などの代数構造と対応づける研究方針に繋がります。
\end{enumerate}

\begin{center}\rule{0.5\linewidth}{0.5pt}\end{center}

\subsection{第4段階：背景時空の排除と自己組織化の力学}\label{ux7b2c4ux6bb5ux968eux80ccux666fux6642ux7a7aux306eux6392ux9664ux3068ux81eaux5df1ux7d44ux7e54ux5316ux306eux529bux5b66}

\textbf{【木原】}
四次元背景時空を仮定しないで、有界なヒルベルト空間から閉じた四次元時空と十分な内部自由度を持つ系が3次元でも10次元でもなく四次元になる理由を説明できるモデルが、これ以外に見当たりません。

\textbf{【Gemini】}
「背景時空を仮定せず、有界なヒルベルト空間から閉じた4次元時空と十分な内部自由度が立ち上がる」という枠組みには、次元をこの形に固定する次の数学的構造が対応しています。

\begin{enumerate}
\def\labelenumi{\arabic{enumi}.}
\tightlist
\item
  \textbf{有界なヒルベルト空間が「閉じた4次元」と「3自由度」に分裂する構造（ホップ写像
  {[}X5{]}{[}X6{]}）：}\\
  状態ベクトルのノルムが常に1に規格化された有界なヒルベルト空間は高次元球面を形成し、位相の冗長性を割って観測空間を取り出す写像はホップファイブレーション（ホップ束）として記述されます。ノルム保存（中心投影）を壊さずにこの写像を作れるのは実数・複素数・四元数・八元数の球面に限られます。

  \begin{itemize}
  \tightlist
  \item
    \textbf{四元数（\(S^3\)）の階層：} 写像は
    \(S^1 \hookrightarrow S^3 \to S^2\) {[}X5{]}
    となり、観測空間が2次元球面（時間を入れても3次元時空）、内部自由度も円周
    \(S^1\)（1自由度）に留まるため、質量と電荷を分離する自由度が不足します。
  \item
    \textbf{八元数（\(S^7\)）の階層：} 写像は
    \(S^3 \hookrightarrow S^7 \to S^4\) {[}X6{]}
    となり、写像先の底空間が閉じた4次元球面
    \(S^4\)（ウィック回転前のド・ジッター時空に対応する閉じた4次元時空）、ファイバーが3次元球面
    \(S^3\)（\((q_1, q_2, q_3)\) の3自由度）として分離します。
  \end{itemize}
\item
  \textbf{10次元等の余剰次元が不要な理由（フルヴィッツの定理
  {[}X7{]}）：}\\
  10次元や11次元の背景時空を先に仮定する枠組みとは異なり、「有界なヒルベルト空間におけるノルム保存（中心投影の成立）」という要請だけで上限が定まります。フルヴィッツの定理
  {[}X7{]}
  により8次元の次の16次元ではノルム保存が破綻するため、中心投影が成立しません。したがって、「十分な内部自由度を持ちつつ閉じた時空へ写像できる最大の器」を選んだ時点で、閉じた4次元時空（\(S^4\)）＋3次元内部自由度（\(S^3\)）＋固有時（\(\tau\)）という次元構成が候補として定まります。
\item
  \textbf{背景時空を仮定しない意味：}\\
  4次元の入れ物を先に置くのではなく、有界なヒルベルト空間とその代数的な共役構造からの写像として4次元時空が創発すると捉えることで、4次元という次元数を「前提」ではなく「帰結」として扱う方針になります。\((q_1, q_2, q_3)\)
  の具体的な物理ラベルの同定は、\(S^3 \hookrightarrow S^7 \to S^4\)
  の対称性から導かれる後続課題となります。
\end{enumerate}

\textbf{【木原】}
なお、ここで「無限のヒルベルト空間」ではなく「有界なヒルベルト空間」を置いているのは、第2段階で曲率半径
\(R\) の中心投影（事象の地平面）を導入した時点で
{[}S1{]}{[}S2{]}、系が無限に発散しない有界で閉じた球面空間であることが自動的に要請されるためです。

\textbf{【Gemini】} 曲率半径 \(R\)
が全ノルムの上限（閉じた球面幾何）を与えるため、無限領域を前提とする非有界な空間ではなく、コンパクトな有界空間として扱う位置づけになります（さらに以下で状態数
\(N\)
を有限とする場合は、有界性に加えて有限自由度の仮定を置く整理になります）。

また、有界・閉空間に加えて状態数 \(N\)
を有限と仮定する動機は、導出上の次の利点にあります。 1.
\textbf{発散の回避：}
無限遠の境界条件や紫外・赤外発散、無限自由度に伴う測度・表現の非同値性の問題が生じず、すべて有限和として扱える。
2. \textbf{有限代数への帰着：} \(N^2\)
全結合からの因子分解や凝縮体の形成を、連続スペクトルの関数解析ではなく、有限次元の線形代数・八元数の不変部分空間分解（ピアース分解や結合子
\([X, Y, Z] = 0\)
の有限連立条件）および数値シミュレーション可能な問題に落とし込める。 3.
\textbf{総自由度の上限：} 全ノルム（曲率半径 \(R\)）と状態数 \(N\)
の双方が有限であるため、凝縮体と相殺された背景（真空）への自由度配分が閉じる。

（数学的定義に関する注記：標準的なヒルベルト空間は複素数体
\(\mathbb{C}\) 上で定義されるため、非結合的な八元数 \(\mathbb{O}\)
上に直接通常のヒルベルト空間の公理系を置くことはできない。本稿でいう「有界なヒルベルト空間」とは、実次元8（複素4次元
\(\mathbb{C}^4\)、規格化球面
\(S^7\)）の状態空間を基礎とし、その状態間の相互作用・積構造として
\(\mathbb{C}^4 \cong \mathbb{O}\)
の同型を介した八元数の非結合代数を付与する立場をとる。）

\textbf{【木原】}
問題は、四元数・八元数の干渉が「読み出し」である点は整理できたものの、相互作用として自身の内部位相をどう変更するのか（動力学）で3カ月以上行き詰まっている点です
{[}S9{]}{[}S10{]}{[}S11{]}。元々この行き詰まりから四元数・八元数を見出しました。動力学を解くのに必要な自由度が足らず、外から粒子の種類や内部状態をヒルベルト空間とは別に手で置く必要があり、それを粒子自身が保持（レジスト）できる根拠が見出せなかったからです。

\textbf{【Gemini】}
「干渉＝読み出し（射影）」までは、有界なヒルベルト空間（\(S^7\)）から閉じた4次元時空（\(S^4\)）へのホップ写像
{[}X6{]}
という静的な幾何学で整理できます。ところが「相互作用＝内部位相の変更（動力学）」に進むと、通常の量子力学の作法では状態
\(\lvert\psi\rangle\) とは別に時間発展演算子（ハミルトニアン
\(\hat{H}\)）や粒子種・結合定数を外から手で与える必要が生じ、背景時空を消したにもかかわらずルールの入れ物を手作業で持ち込む構造になります。

外からパラメーターを足さずに動力学を構成する候補として、八元数
\(\mathbb{O}\) が持つ次の3つの代数構造を整理しました。

\begin{enumerate}
\def\labelenumi{\arabic{enumi}.}
\item
  \textbf{「状態」と「相互作用」を分離しない：トライアリティ（三対性）{[}X8{]}}\\
  通常のヒルベルト空間では状態ベクトルに外部の作用素（行列）を掛けて位相を変えますが、8次元（八元数）には
  \(\mathrm{Spin}(8)\)
  のトライアリティ（三対性）があります。これにより、「粒子Aの状態（左スピノル）」「粒子Bの状態（右スピノル）」「それらを変換する相互作用（ベクトル）」の3つを同じ8次元空間の元として対等に入れ替えられます。外からハミルトニアンを用意せず、状態
  \(X\) と状態 \(Y\) の積 \(X \cdot Y = Z\)
  自体が内部位相の書き換え（動力学）として機能し、状態そのものが相手にとっての相互作用演算子となる構成が候補になります。
\item
  \textbf{内部位相をねじ曲げる構造：非結合性（結合子）}\\
  四元数 \(\mathbb{H}\) までは結合法則 \((AB)C = A(BC)\)
  が成り立つため、積を重ねても座標系の回転（読み出し方の変更）に留まります。八元数では結合法則が破れ、3つの状態が関わったときに結合子（Associator）
  \[[X, Y, Z] = (XY)Z - X(YZ) \neq 0\]
  が生じます。単独で飛ぶ2元の積では結合法則が保たれて測地線運動に対応し、3元以上の干渉が生じた瞬間に結合子が非ゼロとなって、内部空間
  \(S^3\) \((q_1, q_2, q_3)\)
  上の位相をねじ曲げるトルク（相互作用・ゲージ場に対応する曲率）を発生させる仮説になります。
\item
  \textbf{粒子が自身の種類を保持（レジスト）する根拠：極小イデアルと位相ロック
  {[}X9{]}}\\
  八元数の左乗算の連鎖では、どんな八元数を外から掛けても部分空間の外に出ない不変領域（極小左イデアル）{[}X9{]}
  が現れます。相互作用によって連続的な内部位相（\(S^3\)
  上の位置）が変化しても、「どの極小イデアルに属するか」という離散的枠組み（電荷や色荷に対応する量子数・粒子種）が代数的に保たれる機構が候補になります。幾何学的には、\(S^7\)
  上の相互作用ループで生じる巻き付き数（ホロノミー／ベリー位相）が、連続揺らぎで消えないレジスタとして粒子種と固有時
  \(\tau\) の蓄積を保持する見方になります。
\end{enumerate}

この枠組みにおいて、行き詰まりの焦点は「2つの状態が干渉したときにどの演算規則で次の位相へ遷移するか（時間発展の生成子・更新則）」と「遷移後の状態が特定の粒子種・質量として安定に保たれる機構」のいずれにあるかを整理しました。

\textbf{【木原】}
その構造自体はすでに着想しています。問題は、有界で有限な \(N\)
個の状態から \(N^2\)
の相互作用によってヒルベルト空間全体の更新が起こり続けるはずなのに、なぜ一様にならず、局所的な凝縮体の衝突・反射・分裂・合体・放射・吸収に見える「作用の自由度」と「凝縮体の区別」が生まれるのかが理解できていない点です。

\textbf{【Gemini】} \(N\) 個の八元数状態 \(\{O_1, O_2, \dots, O_N\}\) が
\(N^2\)
の全結合で更新され続けるにもかかわらず、「独立した凝縮体（粒子）」と「局所的な6つの相互作用」へ自己組織化するための作業仮説を、以下の2つの柱で整理しました（具体的更新則と厳密な証明は未導出）。

\paragraph{\texorpdfstring{1. \(N^2\)
の全結合から「局所性」と「凝縮体の区別」が生まれる機構（作業仮説）}{1. N\^{}2 の全結合から「局所性」と「凝縮体の区別」が生まれる機構（作業仮説）}}\label{n2-ux306eux5168ux7d50ux5408ux304bux3089ux5c40ux6240ux6027ux3068ux51ddux7e2eux4f53ux306eux533aux5225ux304cux751fux307eux308cux308bux6a5fux69cbux4f5cux696dux4eeeux8aac}

\begin{enumerate}
\def\labelenumi{\arabic{enumi}.}
\tightlist
\item
  \textbf{破壊的干渉による背景平滑化とコヒーレントな凝縮体の抽出：}\\
  全結合で更新されるペアの大部分（\(O(N^2)\) 個）は内部位相（\(S^3\)
  ファイバー）が無相関です。ホップ写像 \(S^7 \to S^4\)
  で4次元時空へ射影すると、無相関な干渉項は互いに打ち消し合い、一様な曲率半径
  \(R\)
  を持つド・ジッター背景時空（真空の期待値）として平滑化される一方、位相整合（コヒーレンス）を起こしてクラスター化した部分集合だけが相殺を免れ、射影面
  \(S^4\) 上の局所的な曲率集中（固有時 \(\tau\)
  とエネルギーの塊＝凝縮体）として残る、という仮説です。
\item
  \textbf{結合子サイクルの閉包による凝縮体の個体性（区別）：}\\
  四元数 \(\mathbb{H}\)
  までは結合的であるため波が線形に拡散しますが、八元数 \(\mathbb{O}\)
  の結合子 \([O_i, O_j, O_k] = (O_i O_j)O_k - O_i(O_j O_k) \neq 0\)
  を用いると、クラスター内部では非結合的なトルクが循環して自己束縛を作りつつ、クラスター外部の状態
  \(O_{\mathrm{ext}}\)
  に対しては結合子がゼロ（\([O_A, O_B, O_{\mathrm{ext}}] \approx 0\)）となる閉じたイデアル構造（色中性・ゲージ不変に相当）が境界線となり、全結合のヒルベルト空間が独立した局所クラスター（凝縮体）の直積へ因子分解される条件を仮定します。
\item
  \textbf{中心投影による距離フィルタリング：}\\
  \(S^7\) から \(S^4\)
  への中心投影（ベルトラミ型射影）の計量構造上、超球面上で位相角が近い状態同士の干渉は投影面上で急峻な位相勾配（強い相互作用）を生む一方、位相角が離れた状態同士の干渉は地平面
  \(R\)
  付近へ引き伸ばされて勾配が潰れます。これにより、実効的な相互作用が
  \(N^2\) から投影面上で近接した凝縮体同士（\(O(N)\)
  の近接作用）へフィルタリングされる見通しになります。
\end{enumerate}

\paragraph{2.
6つの作用の自由度の対応候補}\label{ux3064ux306eux4f5cux7528ux306eux81eaux7531ux5ea6ux306eux5bfeux5fdcux5019ux88dc}

凝縮体（閉じた結合子ループを持つ八元数クラスター）同士が投影面 \(S^4\)
上で接近し位相が重なったとき、\(N^2\)
の更新則の中で結合子サイクルのトポロジカルな組み替えが起きるとして、次の対応候補を整理しました。

\begin{itemize}
\tightlist
\item
  \textbf{放射（Emission）と吸収（Absorption）--- 四元数部分代数
  \(\mathbb{H}\) の切り離しと取り込み：}\\
  凝縮体（八元数
  \(O = a + bl\)、\(\tau > 0\)）が過剰な内部位相のねじれを抱えた際、余剰分が結合法則を満たす純粋な四元数部分代数
  \(\mathbb{H}\)
  のモードとして切り離されると、非結合性による自己束縛（\(\tau\)）を失うため光円錐上の測地線に沿って離脱する（放射）。逆に、飛来した四元数モード（光子）が凝縮体の
  \(S^3\)
  位相と共鳴して結合子ループに取り込まれると、結合子の循環エネルギー（静止質量・固有時
  \(\tau\) の増加）に変換される（吸収）。
\item
  \textbf{衝突（Collision）と反射（Reflection）---
  ホップ不変量による排他：}\\
  2つの凝縮体 \(A, B\) が接近した際、双方の内部で閉じている
  \(S^3 \to S^4\)
  の巻き付き数（ホップ不変量・トポロジカルチャージ）が融合を許さない場合、間に強い結合子反発トルク
  \([O_A, O_B, O_{\mathrm{vac}}]\)
  が生じて運動量ベクトルが反転し、個体性を保ったまま弾き返される操作に対応する。
\item
  \textbf{合体（Fusion）と分裂（Fission）--- 結合子ループの再編成：}\\
  相対運動エネルギーが非結合性の障壁を超えると、\(A\) と \(B\)
  の結合子ループが \(S^7\)
  上で繋ぎ替わって単一のより大きな閉じたサイクルになる（合体）。逆に、単一凝縮体内部で
  \(q_1, q_2, q_3\)
  の位相不均衡が増大し有界空間の局所ノルム上限（曲率限界）に達すると、トポロジカル不変量の総和を保ったまま2つの安定ループへ分岐する（分裂）。
\end{itemize}

\paragraph{次の焦点}\label{ux6b21ux306eux7126ux70b9}

この幾何学的・代数的枠組みを具体的な動力学方程式（更新則）として書き下す際の急所は、\(N\)
個の八元数状態の更新演算子をノルム保存（ユニタリ性）を保ったまま非結合な積でどう定義するか（「八元数行列の作用」のような線形演算に近い形式か、\(O_i \mapsto O_i + \sum_j [O_i, O_j, \dots]\)
のような状態同士の直接積による非線形な微分・差分方程式か）にあります。

\textbf{【木原】}
つまり八元数は、作用の自由度や凝縮体の因数分解を外から与えずとも定式化できる代数構造を持っている可能性がある、ということですね。

\textbf{【Gemini】}
複素数や四元数といった結合法則が成り立つ数体系では、全結合の系は線形に混ざり合うだけで、「どこで区切るか（因数分解）」や「どう力を及ぼすか（相互作用）」を外から定義する必要があります。一方、八元数では代数構造そのものが次の機構を与える候補になります。

\begin{enumerate}
\def\labelenumi{\arabic{enumi}.}
\tightlist
\item
  \textbf{結合子 \([X, Y, Z]\) による作用の自由度の生成：}
  外からハミルトニアン（力）を与えなくても、3つ以上の状態が絡み合った際に生じる非結合性のズレそのものが、内部空間のねじれ・束縛エネルギー・散乱トルクの候補になります。
\item
  \textbf{非結合的な閉包（イデアル）による凝縮体の因数分解：}
  外部に対して結合子がゼロになる閉じた部分空間へと状態が相分離する条件が成り立てば、全結合の
  \(N^2\)
  から独立した個体（凝縮体）と結合法則を満たす四元数モード（放射・吸収される光子）が外部境界条件なしで切り出されます。
\end{enumerate}

これを具体的な計算・定式化へ落とし込む際の計算上の鍵として、次の2点が挙げられます。

\begin{itemize}
\tightlist
\item
  \textbf{ピアース分解（べき等元による自動因数分解）：}
  八元数（あるいはその複素化・行列拡大であるアルバート代数）の中で
  \(P^2 = P\)
  を満たす原始べき等元を基準に展開すると、外から座標軸や粒子種を設定せずとも空間を時空ブロックと3つの内部自由度へ直交分解する枠組みが得られます。
\item
  \textbf{「交換子＋結合子」による時間発展：} 外部のエネルギー演算子
  \(H\) を用意する代わりに、状態同士の交換子
  \([O_i, O_j]\)（光円錐上の回転・並進）と結合子
  \([O_i, O_j, O_k]\)（凝縮体同士の衝突・合体・分裂）だけで1ステップの更新則（非結合なリウヴィル方程式／マルツェフ力学）を構成し、系の全ノルムを保ったまま自己組織化と因数分解を行う方針になります。
\end{itemize}

\begin{center}\rule{0.5\linewidth}{0.5pt}\end{center}

\subsection{補遺1：量子ビット（ブロッホ球）表現との数学的同型性（論文整理時の追記）}\label{ux88dcux907a1ux91cfux5b50ux30d3ux30c3ux30c8ux30d6ux30edux30c3ux30dbux7403ux8868ux73feux3068ux306eux6570ux5b66ux7684ux540cux578bux6027ux8ad6ux6587ux6574ux7406ux6642ux306eux8ffdux8a18}

\textbf{【木原】} 第3段階に現れたホップ写像 \(S^3 \to S^2\) {[}X5{]}
について、\((z_1, z_2)\) から 3次元座標 \((x, y, z)\)
を構成する式は、量子ビットの表現とよく似ているように見えます。

\textbf{【Gemini】}
おっしゃる通り、それは\textbf{量子ビット（qubit）のブロッホ球（Bloch
sphere）表現と数学的に同一の写像}です。

1量子ビットの純粋状態は2つの複素振幅 \((z_1, z_2) \in \mathbb{C}^2\)
を用いて
\[\lvert\psi\rangle = z_1 \lvert 0\rangle + z_2 \lvert 1\rangle \quad (|z_1|^2 + |z_2|^2 = 1)\]
と表され、状態空間は \(S^3\) です。パウリ行列
\((\sigma_x, \sigma_y, \sigma_z)\)
の期待値（ブロッホベクトル）を計算すると、
\[x = \langle\psi\rvert \sigma_x \lvert\psi\rangle = 2\,\mathrm{Re}(z_1 \bar{z}_2),\quad y = \langle\psi\rvert \sigma_y \lvert\psi\rangle = 2\,\mathrm{Im}(\bar{z}_1 z_2) = -2\,\mathrm{Im}(z_1 \bar{z}_2),\quad z = \langle\psi\rvert \sigma_z \lvert\psi\rangle = |z_1|^2 - |z_2|^2\]
となり、全体のグローバル位相 \(e^{i\phi}\)（\(S^1\)
ファイバー）を商空間で割ってブロッホ球面 \(S^2\)
へ写す操作そのものがホップ写像 \(S^1 \hookrightarrow S^3 \to S^2\)
に対応します（光学における光子の偏光状態とポアンカレ球／ストークスパラメータの対応も同じ構造です）。さらに、2量子ビット（複素4成分
\(\mathbb{C}^4\)、規格化で \(S^7\)）に対して次の階層のホップ写像
\(S^3 \hookrightarrow S^7 \to S^4\) {[}X6{]}
が現れることも知られており、有界なヒルベルト空間から時空・内部空間が分離する本稿の視点と量子状態の幾何学が直結しています。

\begin{center}\rule{0.5\linewidth}{0.5pt}\end{center}

\subsection{補遺2：倍音（局在波・倍音循環）と凝縮体・質量階層への接続課題（論文整理時の追記）}\label{ux88dcux907a2ux500dux97f3ux5c40ux5728ux6ce2ux500dux97f3ux5faaux74b0ux3068ux51ddux7e2eux4f53ux8ceaux91cfux968eux5c64ux3078ux306eux63a5ux7d9aux8ab2ux984cux8ad6ux6587ux6574ux7406ux6642ux306eux8ffdux8a18}

\textbf{【木原】}
初期の研究から、局所性を与えるためには倍音（奇数倍音和など）が必要であることを検討してきました
{[}S12{]}--{[}S15{]}。本稿にはまだ倍音をどのように扱うかという観点が抜けていますが、これは第4段階の「凝縮体」と接続する可能性があり、また八元数（8自由度）だけでは自由度が不足する可能性のある質量階層や粒子の種類を、倍音構造によって説明できる可能性があるのではないでしょうか。

\textbf{【Gemini】}
ご指摘の通り、本稿の八元数・凝縮体モデルと、既発表の倍音シリーズ
{[}S12{]}--{[}S15{]}
を接続することは、理論の欠落を補う上で重要な論点になります。

\begin{enumerate}
\def\labelenumi{\arabic{enumi}.}
\tightlist
\item
  \textbf{局所性の生成と凝縮体のプロファイル：}
  半波長位相区間における奇数倍音和の孤立ピーク波とその局在性
  {[}S12{]}、および局在奇数倍音波の再生核性 {[}S13{]}
  は、第4段階で述べた「相殺を免れて局所的に浮かび上がる凝縮体」の空間局在構造そのものを与える候補になります。
\item
  \textbf{閉鎖系の偶奇倍音と倍音循環による持続性：}
  有限完全関係閉系における偶奇倍音の解析 {[}S14{]}
  や、局在波の倍音交換・倍音循環によって質量獲得と粒子寿命（散逸せず持続する機構）を扱う視点
  {[}S15{]}
  は、凝縮体がなぜ安定に保たれるかを動力学側から支える枠組みと対応します。
\item
  \textbf{八元数の8自由度と質量階層・粒子種の補完：}
  八元数の8自由度（および極小左イデアルによる離散枠組み）が時空と基本対称性の骨格を定める一方、異なる質量階層や多様な粒子種・励起スペクトルを表現するには、単一の8次元元だけでなく「どの倍音次数・倍音モードの組み合わせで凝縮体が構成されているか」というスペクトル自由度が必要になる可能性があります。
\end{enumerate}

したがって、八元数による\textbf{代数的因子分解（個体と対称性の枠組み）}と、倍音による\textbf{局在化・スペクトル構造（凝縮体の内部構造と質量階層）}をどう統合するかが、第4段階に続く重要な接続課題となります。また、第4段階で「状態数
\(N\) が有限」と置きつつ射影先として連続的な球面 \(S^4\)
上の局所凝縮体を扱う際の橋渡しとして、有限個 \(N\)
の離散モード（倍音次数）の重ね合わせが投影面上で局所的な孤立ピーク波（凝縮体）{[}S12{]}
を形成する、という対応づけも今後の検討対象となります。

\begin{center}\rule{0.5\linewidth}{0.5pt}\end{center}

\subsection{参考文献}\label{ux53c2ux8003ux6587ux732e}

\subsubsection{自己引用（Concept
DOI、10.5281/zenodo.〔番号〕）}\label{ux81eaux5df1ux5f15ux7528concept-doi10.5281zenodo.ux756aux53f7}

{[}S1{]} 木原範昭「中心投影による4次元空間の幾何学的定式化」2026-04-05.
Concept DOI 10.5281/zenodo.19427779.

{[}S2{]} 木原範昭「中心投影フレームワークにおける Schwarzschild--de
Sitter 厳密解」2026-04-12. Concept DOI 10.5281/zenodo.19538097.

{[}S3{]} 木原範昭「複素数の集合から時空 (1＋3) への写像 ―
四元数による読み出しと部分投射 ―」2026-10-18. Concept DOI
10.5281/zenodo.23451494.

{[}S4{]}
木原範昭「四元数の関係読み出しと、結晶化としてのインフレーション ―
無名性を保った自律的構造化と部分投射 ―」2026-10-21. Concept DOI
10.5281/zenodo.23523237.

{[}S5{]} 木原範昭「複素数を正しく理解していたか ―
振幅はどこに住むか、虚数読み、\(\mathbb{C}\) → \(\mathbb{H}\) →
\(\mathbb{O}\) の帰属 ―」2026-10-23. Concept DOI
10.5281/zenodo.23585436.

{[}S6{]} 木原範昭「6次元超直方体を採用した (R, Q)
マッピングの構成」2026-04-25. Concept DOI 10.5281/zenodo.19692852.

{[}S7{]} 木原範昭「統合と 4+1 から 3+1
への削減を主要未解決問題として（論文6：統合）」2026-04-28. Concept DOI
10.5281/zenodo.19837597.

{[}S8{]}
木原範昭「無名部分投射による物理的対称性選択の公理的枠組み」2026-07-24.
Concept DOI 10.5281/zenodo.21521899.

{[}S9{]}
木原範昭「無名二チャネル閉鎖波系におけるロック動力学」2026-08-02.
Concept DOI 10.5281/zenodo.21763999.

{[}S10{]}
木原範昭「インフレーションの終わり方が、物質生成の始まり方である------タイミング則を置かない万能相互作用のN体埋め込みと通し創世走行」2026-08-05.
Concept DOI 10.5281/zenodo.21809814.

{[}S11{]}
木原範昭「空間軸3軸と固有時間の創生------凝縮体からの読出しの再導出と、唯一読めない座標時間」2026-08-06.
Concept DOI 10.5281/zenodo.21816651.

{[}S12{]}
木原範昭「半波長位相区間における一定振幅奇数倍音和の孤立ピーク波とその局在性に関する観察」2026-06-30.
Concept DOI 10.5281/zenodo.20833096.

{[}S13{]} 木原範昭「局在奇数倍音波の再生核性によるボルン分布の形の導出 ―
残す公準を二乗則とランダム性に絞る」2026-06-27. Concept DOI
10.5281/zenodo.20965526.

{[}S14{]} 木原範昭「時間発展を頂点化した有限完全関係閉系」2026-09-15.
Concept DOI 10.5281/zenodo.22763329.

{[}S15{]}
木原範昭「第十三思考実験：質量を獲得した粒子は、なぜ散逸して消滅しないのか？」2026-10-10.
Concept DOI 10.5281/zenodo.23266250.

\subsubsection{外部文献}\label{ux5916ux90e8ux6587ux732e}

{[}X1{]} E. Beltrami, ``Saggio di interpretazione della geometria
non-euclidea,'' \emph{Giornale di Matematiche} 6, 284--312 (1868).

{[}X2{]} W. de Sitter, ``On the relativity of inertia. Remarks
concerning Einstein's latest hypothesis,'' \emph{Proc. Kon. Ned. Akad.
Wet.} 19, 1217--1225 (1917).

{[}X3{]} J. C. Maxwell, \emph{A Treatise on Electricity and Magnetism}
(Clarendon Press, Oxford, 1873).

{[}X4{]} L. E. Dickson, ``On Quaternions and Their Generalization and
the History of the Eight Square Theorem,'' \emph{Ann. Math.} 20,
155--171 (1919). DOI 10.2307/1967865

{[}X5{]} H. Hopf, ``Über die Abbildungen der dreidimensionalen Sphäre
auf die Kugelfläche,'' \emph{Math. Ann.} 104, 637--665 (1931). DOI
10.1007/BF01457962

{[}X6{]} H. Hopf, ``Über die Abbildungen von Sphären auf Sphären
niedrigerer Dimension,'' \emph{Fund. Math.} 25, 427--440 (1935).

{[}X7{]} A. Hurwitz, ``Über die Composition der quadratischen Formen von
beliebig vielen Variabeln,'' \emph{Nachr. Ges. Wiss. Göttingen},
309--316 (1898).

{[}X8{]} J. C. Baez, ``The Octonions,'' \emph{Bull. Amer. Math. Soc.}
39, 145--205 (2002). arXiv:math/0105155

{[}X9{]} C. Furey, ``Standard model physics from an algebra?''
Ph.D.~thesis, University of Waterloo (2016). arXiv:1611.09182

\begin{center}\rule{0.5\linewidth}{0.5pt}\end{center}

\subsection{区分}\label{ux533aux5206}

{\def\LTcaptype{none} % do not increment counter
\begin{longtable}[]{@{}
  >{\raggedright\arraybackslash}p{(\linewidth - 2\tabcolsep) * \real{0.5000}}
  >{\raggedright\arraybackslash}p{(\linewidth - 2\tabcolsep) * \real{0.5000}}@{}}
\toprule\noalign{}
\begin{minipage}[b]{\linewidth}\raggedright
区分
\end{minipage} & \begin{minipage}[b]{\linewidth}\raggedright
内容
\end{minipage} \\
\midrule\noalign{}
\endhead
\bottomrule\noalign{}
\endlastfoot
恒等式・定理 & 第1段階 \(ds^2 = -(c dt)^2 + dx^2 + dy^2 + dz^2\)（光子で
\(ds=0\)）、\((x^2+y^2+z^2)+\tau^2=t^2\)／第3段階
ケイリー＝ディクソン構成 {[}X4{]}、フルヴィッツの定理（実ノルム多元体は
1, 2, 4, 8 次元に限る）{[}X7{]}、ホップ写像 \(S^3 \to S^2\) {[}X5{]}
および \(S^7 \to S^4\) {[}X6{]} \\
読み・解釈 & 第1段階
光子の四次元時空における静止、5次元光円錐からの断面／第2段階
\(R^2=t^2+(i\tau)^2\) による中心投影 {[}X1{]}・ド・ジッター時空 {[}X2{]}
との対比、および主観空間の内部観測者にとっての無区別性（ローレンツ符号は物理的前提）{[}S1{]}／第3段階
複素数2自由度から四元数4自由度（3次元方向保持）および八元数8自由度（時空4＋内部4）への自由度要請
{[}S3{]}--{[}S8{]}／第4段階 曲率半径 \(R\)
の導入に伴う空間の有界・閉鎖性（コンパクト球面）の要請
{[}S1{]}{[}S2{]} \\
仮説 & 第2段階 物質による局所曲率半径 \(R\) の縮小としての重力／第4段階
背景時空を置かない有界ヒルベルト空間とノルム保存からの4次元時空・内部自由度の分離、八元数のトライアリティ
{[}X8{]}・非結合性・極小イデアル
{[}X9{]}・ピアース分解による自己組織化と凝縮体の因子分解 \\
未解決 & 読出し（干渉）ではなく内部位相を書き換える動力学の具体的更新則
{[}S9{]}--{[}S11{]}／\(N^2\)
全結合から独立した局所凝縮体へ因子分解される数学的証明／倍音（局在波・倍音循環）{[}S12{]}--{[}S15{]}
と八元数凝縮体・質量階層・粒子種の統合 \\
\end{longtable}
}

\begin{center}\rule{0.5\linewidth}{0.5pt}\end{center}

\subsection{版の履歴}\label{ux7248ux306eux5c65ux6b74}

\begin{itemize}
\tightlist
\item
  v1.0（2026-11-16）：Zenodo
  公開用書式（題名・副題・著者・ORCID・ライセンス・要旨・はじめに・補遺1・補遺2・区分・参考文献・版の履歴）を整備し、自己引用（\texttt{{[}S1{]}}〜\texttt{{[}S15{]}}）および外部文献（\texttt{{[}X1{]}}〜\texttt{{[}X9{]}}）を整理。
\item
  整理・編集環境：Gemini 4 Hydrogen（思考モード）を利用。
\end{itemize}

\end{document}
